\cotefile{157-5}
\scanpage{213}
\mspage{129}
% archive sheet number: 212
% element: text | top, line 1 | Si on a
Si on a
% element: display | top, line 2 | x is colimit over I
\[
x\simeq\varinjlim_I x_i,\qquad \operatorname{card}\mathrm{Fl}\,I\leq\pi,
\]
% element: text | upper third, lines 3-7 | soit I prime
soit $I'$ l'ensemble des \uline{\unsure{sous-diagrammes}} $J\subset I$ finis de $I$, ordonné par l'inclusion, alors $I'$ est un ens ordonné filtrant, et du fait que $C$ est stable par $\underline{\lim}$ finies on trouve un foncteur
% element: diagram | upper third, lines 8-9 | J maps to x J
\diagram{\begin{tikzcd}[ampersand replacement=\&]
J\arrow[r,mapsto]\&x_J=\varinjlim_{i\in J}x_i
\end{tikzcd}\qquad :\quad
\begin{tikzcd}[ampersand replacement=\&]
I'\arrow[r]\&C,
\end{tikzcd}}
% element: text | middle, line 10 | et on a
et on a
% element: display | middle, line 11 | x is limit over J in I prime
\[
x\simeq\lim_{J\in I'}x_J.
\]

% element: text | middle, line 12 | D'autre part on a
D'autre part on a
% element: display | middle, lines 13-14 | cardinal I prime finite subsets
\[
\begin{aligned}
\operatorname{card}I'&=\operatorname{card}\unsure{\mathscr P|\mathfrak P|\mathrm P}_f(\underbrace{\mathrm{Fl}\,I)}_{\substack{\text{de card}\\\leq\pi}})\leq\pi+\pi^2+\cdots+\pi^{\unsure{n|3}}\unsure{,|+\cdots|\cdots}\\
&\leq\pi\aleph_0=\pi,
\end{aligned}
\]
% element: text | lower middle, lines 15-16 | qed Pour établir
qed. Pour établir 2\textsuperscript{o}), il reste à établir
% element: diagram | lower third, line 17 | P pi implies c and d
\diagram{\begin{tikzcd}[ampersand replacement=\&]
\&P(\pi)\arrow[dl,Rightarrow]\arrow[dr,Rightarrow]\&\\
c\&\&d,
\end{tikzcd}}
% element: text | lower third, lines 18-23 | et l'implication inverse
et l'\deleted{\unsure{implication}}implication inverse $c\Rightarrow P(\pi)$ quand $\pi\geq\pi_1$. Mais $P(\pi)\Rightarrow c$ est une tautologie, et $P(\pi)\Rightarrow d$ puisque $M_\pi$ est stable par $\underline{\lim}$ de cardinal $\leq\pi$. Reste à prouver que $c\Rightarrow P(\pi)$ si $\pi\geq\pi_1$,
% element: deletion | bottom, last two lines | une fois de plus qu'on ait
\deleted{\unsure{une fois de plus} qu'on ait les $\operatorname{card}\mathrm{Hom}(\unsure{x|f},y)\leq\pi$ pour $x,y\in\mathrm{Ob}\,C$.}

\scanpage{214}
\mspage{130}
% archive sheet number: 213
% element: text | top, lines 1-2 | Supposons donc
Supposons donc qu'on ait une famille épimorphique stricte
% element: diagram | upper third, line 3 | x alpha to x
\diagram{\begin{tikzcd}[ampersand replacement=\&]
x_\alpha\arrow[r]\&x
\end{tikzcd}\qquad $\inserted{\alpha\in A,}\ (\operatorname{card}A\leq\pi)$,}
% element: text | upper third, lines 4-8 | cela signifie donc
cela signifie donc que $x$ est limite inductive \inserted{d'un syst. inductif} \unsure{indexé par|indexé sur} la catégorie \inserted{$B$} ayant comme ens. d'objets $A\amalg(A\times A)$, les flèches \deleted{\unsure{étant}} \inserted{non identiques} représentées par les flèches canoniques
% element: diagram | upper middle, lines 9-10 | fibre product and two projections
\diagram{\begin{tikzcd}[ampersand replacement=\&]
\&x_\alpha\times_x x_\beta\arrow[dl,"p_{\alpha,\beta}"']\arrow[dr,"q_{\alpha,\beta}"]\&\\
x_\alpha\&\&x_\beta.
\end{tikzcd}}

% element: text | middle, lines 11-16 | NB B n'est pas
(\uline{NB} $B$ n'est pas une catégorie ordonnée, car pour l'objet $(\alpha,\alpha)$, représenté par $\unsure{x_\alpha\times_x x_\alpha|x_\alpha\times x_\alpha}$, il y a \uline{deux} flèches \unsure{de|vers} l'objet \deleted{$\unsure{\alpha}$}, représenté par $x_\alpha$.\unsure{)|\}} Mais on a \unsure{certainement} $\operatorname{card}\mathrm{Fl}(B)\leq\pi$) Si on savait que $C$ est stable par \deleted{\unsure{produits}} \inserted{\unsure{produits binaires et}} aux sous-objets (ici les $x_\alpha\times_x x_\beta\subset x_\alpha\unsure{\times|<}x_\beta$), on aurait prouvé la condition a). À défaut, \deleted{\unsure{pour|on}} \unsure{utilisons|utilise} \inserted{donc} le fait que $C$ est génératrice par épim. ($=$ stricts) et \unsure{choisissons} pour \unsure{tt|$\forall$} couple $(\alpha,\beta)$ un \deleted{famille} \inserted{ens} épimorphique de flèches :
% State question: the bar in tt / forall is treated as intrinsic, not a separate cancellation.
% State question: the short word before utilisons is barred/overwritten and treated as cancelled.
% element: diagram | lower third, lines 21-23 | x alpha beta sigma in Ob C zero
\diagram{\begin{tikzcd}[ampersand replacement=\&,row sep=1.5em]
x_{\alpha,\beta,\unsure{\sigma|\gamma}}\arrow[r,"u_{\alpha,\beta,\sigma}"]\arrow[d,phantom,"\in"{description,rotate=-90}]\&x_\alpha\times_x x_\beta,\\
\mathrm{Ob}\,C_0\&
\end{tikzcd}}
% element: text | foot, last three lines | où C zero
où $C_0$ est une sous-cat. réduite exhaustive de $C$. Les flèches $u_{\alpha,\beta,\sigma}$ sont connues quand on connaît les composés

\scanpage{215}
\mspage{131}
% archive sheet number: 214
% element: diagram | top, first two lines | components to x alpha and x beta
\diagram{\begin{tikzcd}[ampersand replacement=\&]
x_{\alpha,\beta,\unsure{\sigma|\gamma}}\arrow[r,Rightarrow,"\unsure{v|\nu}_{\alpha,\beta,\sigma}"]\arrow[dr,dash,"\unsure{w|\omega}_{\alpha,\beta,\sigma}"']\&x_\alpha\\
\&x_\beta,
\end{tikzcd}}
% element: text | upper third, lines 3-9 | donc pour
donc pour $x_{\alpha,\beta,\sigma}\in\mathrm{Ob}\,C_0$ fixés, les choix possibles pour $u_{\alpha,\beta,\sigma}$ sont de cardinal majoré par $\pi_1\cdot\pi_1^{\unsure{2|3}}=\pi_1$, de plus on a $\operatorname{card}\mathrm{Ob}\,C_0\leq\pi_1$ donc le nb \deleted{\unsure{total}} des flèches $u_{\alpha,\beta,\sigma}$ \deleted{$\unsure{v_{\alpha,\beta,\sigma},w_{\alpha,\beta,\sigma}}$ \unsure{est}} \inserted{(pour $\alpha,\beta$ fixés)} est majoré par $\pi_1\cdot\pi_1=\pi_1$. Donc le nb. total des flèches $u_{\alpha,\beta,\sigma}$, pour $\alpha,\beta,\sigma$ variables, est majoré par $\pi\cdot\pi_1=\pi$ (si $\pi\geq\pi_1$).

% element: text | middle, lines 10-12 | Donc x est représenté
Donc $x$ est représenté par une catégorie $I$, ayant \deleted{\unsure{les mêmes}} \inserted{comme} objets \deleted{que $B$} \inserted{\unsure{les élém. de}} \deleted{\unsure{la}} somme de $A$ et de
% element: deletion | middle right, overwritten expression | indexed coproduct
\deleted{$\unsure{\coprod_{\alpha,\beta}\Sigma_{\alpha,\beta}|p_{\alpha,\beta}\Sigma_{\alpha,\beta}|\coprod_{\alpha,\beta}\{\sigma\}}$}
% element: deletion | middle, crossed three-line former description | mais les flèches non identiques
\deleted{\unsure{mais} les flèches non identiques de $B$ \unsure{ét} \unsure{étant|A. et} \unsure{remplacées chacune} par les paquets de flèches $v_{\alpha,\beta,\unsure{\lambda|\sigma}},w_{\alpha,\beta,\lambda}$}
% element: text | lower third, lines 15-19 | des indices triples
des indices triples $\alpha,\beta,\sigma$, et pour chaque tel indice, existent deux flèches $v_{\alpha,\beta,\sigma},w_{\alpha,\beta,\sigma}$ non identiques de source cet \deleted{\unsure{indice}} objet, \deleted{\unsure{les|en}} et de but respectivement en $\alpha,\beta$. On a \deleted{\unsure{en fait}} donc $\operatorname{card}\mathrm{Fl}\,I\leq\pi$, et on a un foncteur
% element: diagram | lower third, line 20 | I to C
\diagram{\begin{tikzcd}[ampersand replacement=\&]
I\arrow[r]\&C
\end{tikzcd}}
% element: text | foot, lines 21-24 | alpha maps to x alpha
$\unsure{\alpha_\beta|\alpha}\unsure{\mathrel{\vert\!-\!\vert}|\mapsto}x_\alpha$, $(\alpha,\beta,\sigma)\mapsto x_{\alpha,\beta,\sigma}$, tel que $x=\lim_{i\in I}\unsure{x'|x^i|x_i}$. Cela achève la dém. \unsure{du|de} 1\textsuperscript{o}).

\scanpage{216}
\mspage{132}
% archive sheet number: 215
% element: text | top, lines 1-8 | Prouvons 2 degrees
\uline{Prouvons 2\textsuperscript{o}}. La première \deleted{\unsure{implication}} inclusion (4.5.13.1) est tautologique, la deuxième ne fait qu'exprimer $P(\pi)\Rightarrow d$. Le fait que $\pi^{\pi_0}=\pi$ implique l'égalité des conditions extrêmes est un rappel de 4.5.6. Reste à prouver que $\mathrm{Filt}'{}^\pi$ est stable dans $M$ par sommes de cardinal $\leq\pi$. Mais
% element: text | middle, lines 9-17 | si x alpha
si $(x_\alpha)_{\alpha\in A}$ est une famille d'objets de $\mathrm{Filt}'{}^\pi$, avec $\operatorname{card}A\leq\pi$, et si $x=\coprod_{\alpha\in A}x_\alpha$, prenons pour $\unsure{\forall|\mathrm{tt}}\alpha$ un foncteur $\unsure{\varphi|\phi}_\alpha:I_\alpha\unsure{\longrightarrow|\mathrel{-}}C$, avec $\operatorname{card}\mathrm{Fl}\,I_\alpha\leq\pi$, tel que $x_\alpha$ soit la $\unsure{\varinjlim|\underline{\lim}}$ dans $M$ de ce foncteur, on trouve que $x$ est $\varinjlim$ du foncteur $\varphi:I\longrightarrow C$ défini par les $\varphi_\alpha$, où $I=\coprod_\alpha I_\alpha$ (somme dans $\mathrm{Cat}$). Il est évident que $\operatorname{card}\mathrm{Fl}(I)\leq\pi\cdot\pi=\pi$, ce qui achève de prouver 2\textsuperscript{o}.

% element: text | foot, lines 18-22 | Prouvons 3 degrees
\uline{Prouvons} 3\textsuperscript{o}. On a les implications tautologiques
% element: deletion | foot, cancelled implication chain | 1 implies 1 prime implies 2
\deleted{$1)\unsure{\Leftrightarrow|\Rightarrow}1')\Rightarrow2$}
\inserted{$1\Rightarrow1'\unsure{)|,}$ et $\unsure{2'|1'}\Rightarrow1$} puisque $\mathrm{Filt}'{}^\pi$ est stable par sommes de cardinal $\leq\pi$. Donc $1\Leftrightarrow2'$, et implique que pour $\forall x$ dans $M$,

\scanpage{217}
\mspage{133}
% archive sheet number: 216
% element: text | top, lines 1-4 | I equals Filt prime over x
$I=\mathrm{Filt}'{}^\pi_{/x}$ est stable par $\underline{\lim}$ de cardinal $\leq\pi$, donc est grd devant $\pi$, comme $x\simeq\varinjlim_{i\in I}x_i$, cela prouve que $1\Rightarrow2$. L'implication $2\Rightarrow3$ est \unsure{standart|standard}, \inserted{et résulte de $\mathrm{Filt}'{}^\pi\subset M_\pi$,} je vais expliciter dans un

% element: text | upper middle, lemma heading and introduction | Lemme 4.5.14
\uline{Lemme} 4.5.14. Soit $M$ une catégorie, \inserted{(stable par petites $\varinjlim$), $\pi$ un card. infini,} $M_0^{\unsure{\sharp|*}}$ une sous-catégorie \inserted{str.} pleine de $M$. \deleted{\unsure{telle|stable}}
% element: deletion | middle, cancelled two-line draft | stable par limites inductives
\deleted{stable \inserted{dans $M$} par limites inductives \unsure{de} cardinal $\leq\pi$, \unsure{et} Alors Soit $M_\pi$}
% element: text | middle, introduction to conditions | Considérons les conditions
Considérons les conditions

% element: text | middle, condition a | M zero sharp generating
a) $M_0^\sharp$ est une sous-catégorie \inserted{strictement} génératrice, \deleted{\unsure{stricte}} et stable \inserted{dans $M$} par $\varinjlim$ de cardinal $\leq\pi$.

% element: text | middle, condition b lead | Pour tt objet de M
b) \inserted{Pour} \unsure{Tout|tt} objet de $M$, \deleted{\unsure{est isomorphe à un}} $\exists$ un
% element: display | middle lower, condition b | x is a colimit
\[
x\simeq\varinjlim_{i\in I}\unsure{x\cdot|x_i},
\]
% element: text | lower middle, condition b end | les x i dans M zero sharp
les $x_i$ dans $M_0^\sharp$, \deleted{\unsure{et est}} $I$ grd devant $\pi$.

% element: text | lower middle, condition b prime | Comme b
b') Comme b, avec $I$ cat. \uline{ordonnée}.

% element: text | lower third, condition c | M pi contained in M zero
c) $M_\pi\subset M_0$ (où $M_\pi$ est \unsure{le|la} sous-cat. str. pl. de $M$ \unsure{formée|fermée} des objets $\pi$-accessibles).

% element: text | lower third, implications introduction | On a donc les \unsure{implications|conditions}
On a donc les \unsure{implications|conditions}
% element: diagram | lower third, implication chain | a implies b equivalent b prime implies c
\diagram{\begin{tikzcd}[ampersand replacement=\&]
a\arrow[r,Rightarrow]\&b\arrow[r,Leftrightarrow]\&b'\arrow[r,Rightarrow]\&c,
\end{tikzcd}}
% element: margin note | lower left slanted beside conclusion | cf dem de prop 4.5.6
% slanted
\margin{cf dém. de (\inserted{\unsure{prop.|prop} 4.5.6,} $4\Rightarrow5$), page 117}
% element: text | foot, conclusion | et si on suppose
et si on suppose $M_0\subset M_\pi$, \inserted{et $M_\pi$ str. gén. de $M$,} alors les quatre conditions a), b, b', c sont équivalentes, et \unsure{équivalentes} à $M_{\unsure{0|\pi}}=M_\pi$.

\scanpage{218}
\mspage{134}
% archive sheet number: 217
% element: text | top, lines 1-3 | Enfin revenant a la dem
\deleted{La dém.} Enfin (revenant à la dém. de 4.5.13, 3\textsuperscript{o}), l'implication $\unsure{3|B})\Rightarrow1)$ provient du fait que $M_\pi$ est stable par $\varinjlim$ de cardinal $\leq\pi$.

% element: text | upper third, lines 4-9 | Reste a prouver
Reste à prouver que ces conditions, \unsure{ici|i.e.} la stabilité de $\mathrm{Filt}'_\pi$ par composition de autres flèches, est satisfaite dans les cas A) et C), le cas A) étant d'ailleurs déjà connu. Le cas B) résulte \deleted{\unsure{facile|facilement}} immédiatement de la transitivité des familles épimorphiques strictes universelles, \deleted{\unsure{Je crois que le cas}} et c'est d'ailleurs \unsure{écrit en toutes} explicitement \deleted{\unsure{dans}} (la stabilité \inserted{de $\mathrm{Filt}'{}^\pi$} par $\varinjlim$ de cardinal $\leq\pi$) dans SGA 4 I 9.13. (\unsure{Non pas} dans l'énoncé, mais c'est rappelé dans la démonstration).

% element: text | middle lower, recollection of SGA | En redigeant
En rédigeant, \inserted{SGA 4 I 9} je savais donc que $M_\pi\subset\mathrm{Filt}'{}^\pi$ pour $\forall\pi$, mais l'inclusion, plus évidente pour \unsure{tout|tt} \inserted{(dans le cas $\unsure{\pi=\pi_0|\pi\geq\pi_0}$)}, $\mathrm{Filt}'{}^\pi\subset M_\pi$ (quand on suppose les objets de $C$ $\unsure{\pi_0|\pi_p}$-accessibles) m'avait apparemment échappé !

% element: text | lower third, tentative comparison | Reste a traiter le cas C
Reste à traiter le cas C. J'ai l'impression qu'il est \inserted{un peu} plus général que le cas B), i.e. que \inserted{la condition sur $M$} B implique C, mais non inversement. Ainsi la condition C est vérifiée dès que $M$ \unsure{correspond} : une espèce de \unsure{structures|structure} alg. \unsure{discrète|abstraite} \deleted{\unsure{par|pour}} \inserted{un}

\scanpage{219}
\mspage{135}
% archive sheet number: 218
% element: text | top, lines 1-4 | en termes de limits finies
en termes de $\varprojlim$ \uline{finies}, i.e. quand $M$ est de la forme $\mathrm{Ind}(C)$, $C$ stable par $\varinjlim$ \uline{finies}, alors que je doute que la condition B) soit tjs satisfaite pour ces catégories.

% element: text | upper third, lines 5-9 | Par contre dans la categorie Cat
Par contre dans la catégorie $\mathrm{Cat}$ (qui est un \unsure{prototype} de cette espèce), les $\varinjlim$ filtrantes sont exactes à \unsure{g|e}), mais je ne crois pas que les familles épimorphiques \deleted{\unsure{strictes}} soient universelles. Prenons p.ex. deux hom. de groupes
% element: diagram | middle upper, two homomorphisms | G zero to G and G one to G
\diagram{\begin{tikzcd}[ampersand replacement=\&]
G_0\arrow[r,"u_0"]\&G,\&G_1\arrow[r,"u_1"]\&G
\end{tikzcd}}
% element: text | middle, hypothesis lead | tels que
tels que
% element: diagram | middle, free product map | G zero star G one to G
\diagram{\begin{tikzcd}[ampersand replacement=\&]
G_0*G_1\arrow[r]\&G
\end{tikzcd}}
% element: text | middle, assertion about the family | soit surjectif
soit surjectif, alors la famille \inserted{\unsure{à une limite|d'une limite|non évidente}} est \uline{ép. stricte} dans $\mathrm{Gr}$ et $\mathrm{Cat}$, mais n'est pas \deleted{\unsure{épi}} forcément épimorphique stricte universelle, $\unsure{\overline{\mathrm{m}}|\simeq}$ si
% element: margin note | left middle, slanted correction | c est faux
% slanted
\margin{c'est \uline{faux} ! \deleted{\unsure{l'exemple est concluant|l'argument est concluant}} \unsure{Qui est|qui est} mais vrai dans le cas particulier plus bas.}
% element: text | middle lower, example | p ex G zero equals G one
p.ex. $G_0=G_1=H$, $G=H\times H$, et $u_0,u_1$ sont les injections
% element: diagram | lower middle, injections | H to H times H
\diagram{\begin{tikzcd}[ampersand replacement=\&]
H\arrow[r,shift left]\arrow[r,shift right]\&H\times H,
\end{tikzcd}}
% element: text | lower third, injections values | donnes par
\unsure{données|donnés} par $u_0(\unsure{h|u})=(h,1)$ et $u_1(h)=(1,h)$. Si je prends l'application diagonale
% element: diagram | lower third, diagonal plain line | H line H times H
\diagram{\begin{tikzcd}[ampersand replacement=\&]
H\arrow[r,dash]\&H\times H,
\end{tikzcd}}
% element: text | lower third, inverse images | les images inverses
les images inverses de $u_0,u_1$ sont le $\overline{\mathrm{m}}$ objet de $M_{/H}$, savoir
% element: diagram | lower third, identity group inclusion | e to H
\diagram{\begin{tikzcd}[ampersand replacement=\&]
e\arrow[r]\&H,
\end{tikzcd}}
% element: text | lower third, end of example | et cette famille
et cette \deleted{\unsure{même}} famille n'est ép. str. que si $H$ est le groupe unité.]

% element: text | foot, transition | Le cas C va resulter
Le cas C va résulter du résultat plus précis :

% element: text | foot, corollary heading and first lines | Corollaire 4.5.15
\uline{Corollaire} 4.5.15. \unsure{Dans|Sous} les conditions générales de 4.5.13, supposons que dans $M$ les $\varinjlim$

\scanpage{220}
\mspage{136}
% archive sheet number: 219
% element: text | top, lines 1-4 | filtrantes commutent aux chgt de base
\uline{filtrantes} commutent aux \unsure{chgt.|chgts.} de base (\unsure{p.ex.|i.e.} les $\varinjlim$ filtrantes sont exactes). Alors pour \unsure{tt|$\forall$} $\pi\geq\pi_1$, $\mathrm{Filt}'{}^\pi\unsure{\sharp|\natural|\S}$ est \unsure{formée|fermée} des $x\in\mathrm{Ob}\,M$ satisfaisant la condition suivante.

% element: text | upper third, condition c prime | Il existe une famille de fleches
c$'$) Il existe une fam\uline{ill}e de flèches $(u_\alpha)_{\alpha\in A}$, $u_\alpha:x_\alpha\longrightarrow x$, dans $M$, qui est \uline{épimorphique}, \uline{\unsure{stricte|strictes} \unsure{universelles|universelle}} (i.e. couvrante pour la top. canonique de $M$), avec $x_\alpha\in\mathrm{Ob}\,C$ $\forall\alpha\in A$, et $\operatorname{card}A\leq\pi$.

% element: text | middle upper, demonstration | Dem il faut prouver
\uline{Dém} : il faut prouver que cette condition équivaut à la condition c) de 4.5.13 1°). Comme on a $c'\Rightarrow c$ triviale, il reste à montrer que $c\Rightarrow c'$. Mais on a \unsure{vu|su} (\inserted{4.5.13} 1°)) que c implique $\unsure{a|c}''$), i.e. qu'on a un iso
% element: display | middle, presentation of x | x equals colimit x i
\[
x=\varinjlim_{i\in I}\unsure{x'|x^i|x_i},
\]
% element: text | middle lower, indexing and universal assertion | les xi dans C
les $x_i$ dans $C$, et $I$ ens \uline{ordonné} \uline{filtrant}, avec $\operatorname{card}I\leq\pi$. Il reste à voir que si $x$ est \inserted{dans $M$} une limite inductive \uline{filtrante} \inserted{d'objets $x_i$}, alors \unsure{l'ens|l'une} des flèches $x_i\longrightarrow x$ est épim. stricte \uline{universelle}. Or cela est immédiat, car pour $x'\longrightarrow x$, on a
% element: display | lower third, base changed colimit | x prime equals colimit x prime i
\[
x'=\varinjlim x'_i\qquad\text{(par hyp. C)},
\]
% element: text | lower third, pullback family and cancelled syntax | ou x prime i
\unsure{où|ou} $x'_i=x'\times_x x_i$, \deleted{donc} \inserted{la famille des} $x'_i\longrightarrow x'$ \deleted{est} épim. \unsure{strictes|stricte}, \unsure{par.|pr.}

% element: text | foot, next proof | Prouvons maintenant
Prouvons maintenant que \inserted{dans le cas C,} la condition de 3°, \deleted{\unsure{1}} est satisfaite, i.e. que $\mathrm{Filt}'{}^\pi$ est stable par $\varinjlim$ de cardinal $\leq\pi$. Cela

\scanpage{221}
\mspage{137}
% archive sheet number: 220
% element: text | top, continuation | resulte aussitot du
résulte aussitôt du

% element: text | top, corollary heading and rewritten hypothesis | Corollaire 4.5.16
\uline{Corollaire} 4.5.16 (de 4.5.15). Dans les conditions de 4.5.1\deleted{\unsure{3}}5, \deleted{(\unsure{où|ou} \unsure{l'hypothèse|la condition} C) de 4.5.13 \unsure{(3 C)|3 C})} \inserted{(i.e. la condition C) de 4.5.13.3°)},
% element: deletion | upper third, crossed line with dotted strokes beneath | soit x un objet de M
% Short separate strokes run below this crossed line and below existe; unsure alternatives explicitly preserve cancelled versus reinstated state.
\unsure{\deleted{soit $x$ un objet de $M$, et supposons qu'il existe}|soit $x$ un objet de $M$, et supposons qu'il existe} une famille \uline{épimorphique} \uuline{universelle} $(x_\alpha\longrightarrow x)_{\alpha\in A}$ de flèches \inserted{de $M$}, avec $x_\alpha\in\mathrm{Ob}\,\mathrm{Filt}'{}^\pi$ pour $\forall\alpha\in A$, et $\operatorname{card}A\leq\pi$. Alors $x\in\mathrm{Ob}\,\mathrm{Filt}'{}^\pi$.

% element: text | middle upper, proof | En effet l hyp sur x alpha
En effet, l'hyp. \unsure{sur|\deleted{sur}} $x_\alpha$ \deleted{implique} \unsure{équivaut|équivaut à}, en vertu de 4.5.15, qu'il existe une famille \inserted{épim. stricte universelle} \unsure{de} flèches
% element: diagram | middle, covering family of x alpha | x alpha beta indexed by beta in B alpha to x alpha
\diagram{\begin{tikzcd}[ampersand replacement=\&]
(x_{\alpha\beta})_{\beta\in B_\alpha}\arrow[r]\&x_\alpha,
\end{tikzcd}}
% element: text | middle, cardinal bound | avec card B alpha
\deleted{\unsure{pour un}} \deleted{de} \inserted{avec} $\operatorname{card}B_\alpha\leq\pi$ et les $x_{\alpha\beta}$ dans $C$. Soit
% element: display | middle, disjoint union of index sets | B equals coproduct B alpha
\[
B=\coprod_{\alpha\in A}B_\alpha,
\]
% element: text | middle lower, composite family | on a alors card B
on a alors $\operatorname{card}D\leq\pi^2=\unsure{\pi}$, et d'autre part la famille des flèches composées
% element: diagram | lower middle, composite arrows | x alpha beta to x alpha to x
\diagram{\begin{tikzcd}[ampersand replacement=\&]
x_{\alpha\beta}\arrow[r]\&x_\alpha\arrow[r,dash]\&x
\end{tikzcd}}
% element: text | lower third, conclusion | indexee par B
indexée par $B$ est épimorphique stricte universelle (par transitivité de ce type de familles), donc $x\in\mathrm{Ob}\,\mathrm{Filt}'{}^\pi$, qed.

% element: drawing | space above final paragraph | isolated oblique stroke
\figure{\begin{tikzpicture}\draw (0,0)--(0.35,-0.5);\end{tikzpicture}}

% element: text | foot, transition to characterization | Avec tout ca
Avec tout ça, on n'a toujours pas donné de caractérisation de $M_\pi$ lui-même, qui est \unsure{meilleur|meilleure} que $\mathrm{Filt}^\pi$ et $\mathrm{Filt}'{}^\pi$ du point de vue \unsure{catégoriques|catégorique}, \unsure{car} \inserted{seuls} les $M_\pi$ \deleted{\unsure{forment}} définissent une filtration cardinale, et de plus ils ont les propriétés de stabilité désirables dans \unsure{ces|les}

\scanpage{222}
\mspage{138}
% archive sheet number: 221
% element: text | top, continuation and criterion | contextes pour les limites
\unsure{contextes}, pour les $\unsure{\underline{\lim}|\varinjlim}$ de cardinal $\leq\pi$. Le \inserted{seul} critère non \unsure{tautologique|catégorique} que j'arrive à dégager dans le cas général (pour \unsure{tt|tout|$\forall$} $\pi\geq\pi_0$, et sans hyp. \unsure{d'existence|d'exactitude} supplémentaires sur $M$) est le suivant, beaucoup moins ``explicite'' \deleted{\unsure{que|pour}} ou ``calculatoire'' que \deleted{les deux} les critères très concrets pour $\mathrm{Filt}^\pi$ et pour $\mathrm{Filt}'{}^\pi$.

% element: text | middle upper, proposition | Proposition 4.5.17
\uline{Proposition} 4.5.17. Soit $\pi$ un cardinal $\geq\pi_0$. Alors $M_\pi$ est \unsure{égale|égal} à la plus petite sous-catégorie strictement pleine de $M$, contenant $C$, et stable par $\varinjlim$ de cardinal $\leq\pi$.

% element: text | middle, bracketed qualification | Il est important de noter
[Il est important de noter qu'il ne \uline{suffit pas} de dire stable par $\underline{\lim}_I$; \deleted{de cardinal $\leq\unsure{{-}|\pi}$} \inserted{pour $I$ grd devant $\pi_0$} et $\operatorname{card}I\leq\pi$. On \deleted{\unsure{risque de}} \unsure{tombe}\deleted{\unsure{r}} \inserted{dans} \inserted{$\mathrm{Filt}^\pi$, qui risque d'être} une catégorie \deleted{\unsure{strictement}} plus petite que $M_\pi$. ] \inserted{Ce n'est pas non plus la cat. des $\underline{\lim}_I \unsure{x\cdot|x_i}$, $x_i\in C$ $\forall i\in\mathrm{Ob}\,I$, et $\operatorname{card}\unsure{\mathrm{Fl}|\mathbb{F}\mathrm{l}}\,I\leq\pi$ --- c'est \unsure{presque|précisément|puisque} cette catégorie --- c'est $\mathrm{Filt}'{}^\pi$.]}

% element: text | lower third, proof heading and inclusion | Dem M prime
\uline{Dém} (\unsure{soit} $M'$ la cat. de l'énoncé) on a bien $M'\subset M_\pi$, puisque $M_\pi$ contient $C$ et est stable par les $\underline{\lim}$ de cardinal $\leq\pi$. Inversement, le lemme 4.5.14
% element: diagram | foot left, implication reference | a implies c
\diagram{\begin{tikzcd}[ampersand replacement=\&]
a\arrow[r,Rightarrow]\&c:
\end{tikzcd}}
% element: drawing | foot centre | isolated oblique stroke
\figure{\begin{tikzpicture}\draw (0,0)--(0.35,-0.7);\end{tikzpicture}}

\scanpage{223}
\mspage{139}
% archive sheet number: 222
% element: text | top, section heading | 4.6 Caracterisation des cat proto-topos
4.6. \uline{Caractérisation des cat. \unsure{proto-topos|prototopos}} (dém. \unsure{de|du} 4.3.6)

% element: text | upper third, return to earlier proof | Il est temps de revenir
\deleted{4.6.1.} Il est temps de revenir à la démonstration du th. 4.3.6, \inserted{(p. 70--71)} celle que j'avais donnée \unsure{étant|était} incomplète (et erronée --- j'affirmais que $(\mathrm{Filt}^\pi)_{\pi\geq\pi_{\unsure{1|0}}}$ est une filtration cardinale, ce qui est faux si on n'admet que $\pi^{\pi_0}=\pi$ pour $\unsure{\forall|\mathrm{tt}}\pi\geq\pi_1$.

% element: text | middle upper, logical scheme lead | On avait vu l equivalence
On avait vu l'équivalence des conditions a, a$'$, b, soit PTA. Je vais donner la démonstration par le schéma logique
% element: diagram | middle, implication chain | PTA implies d implies c implies PTA
\diagram{\begin{tikzcd}[ampersand replacement=\&]
(\unsure{\alpha|a})\&\mathrm{PTA}\arrow[r,Rightarrow]\&d)\arrow[r,Rightarrow]\&c)\arrow[r,Rightarrow]\&\mathrm{PTA}.
\end{tikzcd}}
% element: text | middle, accessible proto-topos and small category | Prouvons PTA implies d
Prouvons $\mathrm{PTA}\Rightarrow d)$, i.e. que si $M$ est un proto-topos accessible, alors \deleted{\unsure{il y a|il existe}} \inserted{\unsure{existe}} \inserted{un cardinal $\pi_0$ \unsure{infini}, et} une \deleted{\unsure{petite}} \inserted{petite} catégorie $C_0$, \deleted{\unsure{qui|que}} \inserted{stable par $\unsure{\underline{\lim}|\varinjlim|\varprojlim}$ de cardinal $\leq\pi_0$} et une équivalence de catégories
% element: diagram | lower middle, equivalence 4.6.1 | M to Ind prime pi zero C zero
\diagram{\begin{tikzcd}[ampersand replacement=\&]
\text{(4.6.1)}\&M\arrow[r,"\approx"]\&\mathrm{Ind}'_{\pi_0}(C_0).
\end{tikzcd}}
% element: text | lower third, transition | Plus precisement
Plus précisément, on a ceci :

% element: text | lower third, proposition 4.6.1 and corrections | Proposition 4.6.1 Soit M
\uline{Proposition} 4.6.1. Soit $M$ un proto-topos, i.e. une \deleted{proto-topos} \inserted{$\unsure{U|\mathcal{U}|\mathscr{U}}$catégorie \unsure{et} stable par petites $\underline{\lim}$,} admettant une \deleted{petite} sous-catégorie \inserted{pleine} génératrice $S$, \inserted{\unsure{qui est une} petite catégorie} d'objets accessibles de $M$. Soit $\pi_0$ un $\unsure{*|+}$ cardinal tel que les objets de $S$ soient $\pi_0$-\unsure{accessibles|accessible}, \inserted{(i.e. $S\subset M_{\pi_0}$)} \deleted{$\unsure{C|S}$} et soit $C$ la plus petite sous-catégorie str. pleine de $M$, contenant \deleted{$\unsure{C|S}$} $S$ et stable par $\varinjlim$ de cardinal $\leq\pi_0$. Alors
% element: text | foot, first conclusion | C est equivalente
\unsure{1°)|a)} $C$ est équivalente à une petite catégorie.

\scanpage{224}
\mspage{140}
% archive sheet number: 223
% element: text | top, second conclusion | C equals M pi zero
2°) $C=M_{\pi_0}$, sous-catégorie \deleted{\unsure{p|pl.}} str. pleine de $M$ \unsure{formée|formé} des objets $\pi_0$-accessibles.

% element: text | upper third, third conclusion | Le foncteur canonique
3°) Le foncteur canonique
% element: diagram | upper third, canonical functor | Ind prime pi zero C to M
\diagram{\begin{tikzcd}[ampersand replacement=\&]
\mathrm{Ind}'_{\pi_0}(C)\arrow[r,"\approx"]\&M
\end{tikzcd}}
% element: margin note | left upper-middle, slanted | donc aussi
% slanted
\margin{donc \unsure{:|.} \unsure{aussi|ainsi}
% element: diagram | slanted left margin | Ind prime pi C zero to M
\diagram{\begin{tikzcd}[ampersand replacement=\&]
\mathrm{Ind}'_{\pi}(C_0)\arrow[r,"\approx"]\&M
\end{tikzcd}}
où $C_0$ est une \uline{petite} catégorie \unsure{équiv.|équivalente} à $\unsure{C}$.}
% element: text | middle upper, characterization | defini a iso unique pres
(défini à iso \deleted{près} unique près par la condition que son composé avec l'inclusion
% element: diagram | middle, inclusion in Ind | C into Ind prime pi zero C
\diagram{\begin{tikzcd}[ampersand replacement=\&]
C\arrow[r,hook]\&\mathrm{Ind}'_{\pi_0}(C)
\end{tikzcd}}
% element: text | middle, comparison inclusion | soit iso a l inclusion
soit iso à l'inclusion
% element: diagram | middle, inclusion plain stroke | C line M
\diagram{\begin{tikzcd}[ampersand replacement=\&]
C\arrow[r,dash]\&M,
\end{tikzcd}}
% element: text | middle, preservation condition | et que rho commute
et que $\unsure{\rho|\varphi|\phi}$ commute aux \inserted{petites} $\varinjlim$, ou même aux \deleted{petites} $\underline{\lim}$ filtrantes \deleted{de cardinal} \inserted{\unsure{grdes|grds}} devant $\pi_0$ --- cf 4.4.16 \unsure{page|pge} 101) \deleted{est} est une équivalence de catégories.

% element: text | lower half, demonstration | Dem On va construire C
\uline{Dém}. 1°) On va construire $C$ par induction transfinie, en posant
% element: display | lower half, initial stage | C zero generated by S
\[
\begin{aligned}
C_0&=\text{\unsure{sous|Sous}-catégorie str. pleine de $M$ engendrée}\\
&\quad\text{par $S$}
\end{aligned}
\]
% element: display | lower third, successor stage | C alpha plus one
\[
\begin{aligned}
C_{\alpha+1}&=\text{Sous-catégorie str. pleine de $M$ \deleted{\unsure{déf. engendrée|des objets}}}\\
&\quad\text{déf. par les }\varinjlim_I\text{ dans $M$ d'objets de}\\
&\quad C_\alpha,\quad\text{\unsure{où|avec} $I$ catégorie telle que}\\
&\quad\operatorname{card}\mathrm{Fl}\,I\leq\pi_0
\end{aligned}
\]
% element: display | lower third, limit stage | C alpha union earlier C alpha prime
\[
C_\alpha=\bigcup_{\alpha'<\alpha}C_{\alpha'},\quad\text{si $\alpha$ est un ordinal limite.}
\]
% element: text | foot, uniqueness and continuation | Cela definit
% Intersecting earlier and higher strokes at de facon: state/underlying wording remains uncertain.
Cela définit \deleted{\unsure{d'une façon|sans ambiguïté}} \inserted{\unsure{de façon}} unique les $C_\alpha$ (pour $\forall$ petit ordinal $\alpha$). Par la définition de $C$, on voit par induction transfinie que les $C_\alpha$ sont $\subset C$. Soit $\alpha_0$

\scanpage{225}
\mspage{141}
% archive sheet number: 224
% element: text | top, continuation of transfinite construction | le premier ordinal
le premier ordinal tel que $\operatorname{card}\alpha_0>\pi_0$, \inserted{donc l'ens $\bar I_{\alpha_0}$ des ord. $\alpha<\alpha_0$ est grd devant $\pi_0$.} Je dis \deleted{que} qu'on a $C_{\alpha_0}=C$. Pour ceci, il suffit de voir que $C_{\alpha_0}$ est stable par $\unsure{\underline{\lim}|\varinjlim}$ de cardinal $\leq\pi$. Mais si $I$ est une catégorie telle que $\operatorname{card}\mathrm{Fl}\,I\leq\pi_0$,
% element: diagram | upper third, first functor and inserted union | phi I to C alpha zero
\diagram{\begin{tikzcd}[ampersand replacement=\&]
\unsure{\varphi|\phi}:\ I\arrow[r]\&C_{\alpha_0}\inserted{=\bigcup_{\unsure{\alpha'|\alpha^\ast|\alpha^\sharp}<\alpha_0}C_{\unsure{\alpha'|\beta'}}}
\end{tikzcd}}
% element: text | upper middle, factorization | foncteur alors
foncteur, alors la définition de $\alpha_0$ implique que $\varphi$ factorise en
% element: diagram | middle, factorization plain stroke | phi alpha I line C alpha prime
\diagram{\begin{tikzcd}[ampersand replacement=\&]
\varphi_\alpha:\ I\arrow[r,dash]\&C_{\alpha'},
\end{tikzcd}}
% element: margin note | left middle, slanted question and correction | Un peu court
% slanted
% The last two lines are crossed; the large crossing also intersects earlier strokes. The short word after est and the cancelled opening remain uncertain.
\margin{? Un peu court. Il faut invoquer le fait que \deleted{\unsure{card}} $\bar I_{\alpha_0}$ est \unsure{le|b.} grd devant $\pi_0$\unsure{;|,} \deleted{(\unsure{Il|$\bar I_{\alpha_0}$} \unsure{est} \unsure{$=$|$\bar m$} grd devant \unsure{$\forall$|tt} $\unsure{\pi'_0|\pi_0}$-\unsure{quelconque})}}
% element: text | middle, closure conclusion | pour alpha prime less than alpha zero
pour $\alpha'<\alpha_0$. Donc $\underline{\lim}\varphi_{\deleted{\unsure{\alpha_0|\alpha'}}}\in\mathrm{Ob}\,C_{\alpha+1}\subset\mathrm{Ob}\,C_{\alpha_0}$, ce qui montre la stabilité requise.

% element: text | middle lower, essential smallness | Pour voir que C est equiv
Pour voir que $C$ est équiv. à une petite catégorie, il suffit de le voir pour les $C_\alpha$, \unsure{par|en} procédant par induction transfinie. Cela résulte du lemme suivant, immédiat:

% element: text | lower third, lemma 4.6.2 | Lemme La limite inductive
% The initial article is heavily retraced; La, an overwritten Toute, and underline/cancellation states are explicitly alternative readings.
\uline{Lemme} 4.6.\unsure{1|\deleted{1}}2. \unsure{1°)|a)} \unsure{\uline{La}|\deleted{Toute}\uline{La}|\uline{Toute}|\deleted{La}} limite inductive \inserted{$C$} d'un syst. inductif filtrant $(C_i)_{i\in I}$ de catégories $C_i$ équivalentes à des petites catégories, est équivalente à une petite catégorie. Si $\unsure{\operatorname{card\,rél}|\operatorname{card\,réd}}(C_i)\leq\pi$ \inserted{pour $\forall i$,} et $\operatorname{card}I\leq\pi$, alors $\operatorname{card\,red}C\leq\pi$ ($\pi$ étant un

\scanpage{226}
\mspage{142}
% archive sheet number: 225
% element: text | top, continuation and cardinal notation | un cardinal infini
un cardinal infini). Ici, pour la catégorie \deleted{\unsure{la|$\mathcal A$}} $A$, je pose $\operatorname{card}A=\operatorname{card}\mathrm{Fl}(A)$, \deleted{et} $\operatorname{card\,red}(A)=\operatorname{card}(A_0)$, si $A_0$ est une catégorie réduite équivalente à $A$.

% element: text | upper third, second lemma clause | Soit C zero subset M
2°) Soit $C_0\subset M$ une sous-catégorie \inserted{\unsure{pl|de}} d'une catégorie $M$, \inserted{$\pi_0$ un cardinal infini,} et soit $\unsure{C'_0|C_0}$ la sous-catégorie str. pleine de $M$ \deleted{formée} engendrée par les limites inductives \deleted{$\varinjlim$ d'obj} dans $\deleted{\unsure{C|E}}M$ (\unsure{représentables|représentées} dans $\deleted{\unsure{C|E}}M$) des foncteurs
% element: diagram | middle, indexing functor | I to C zero
\diagram{\begin{tikzcd}[ampersand replacement=\&]
I\arrow[r,two heads]\&C_0,
\end{tikzcd}}
% element: text | middle, cardinal bounds | pour I categorie telle que
\unsure{pour|avec} $I$ catégorie telle que $\operatorname{card}\bar I\leq\pi_0$. Si $C_0$ est équivalente à une petite catégorie, il en est de même de $\deleted{\unsure{C|C_0|E}}C'_0$. Plus précisément, $\operatorname{card\,red}(C_0)\leq\pi_1$, \inserted{$\pi_1$ un cardinal infini,} alors \deleted{card} $\operatorname{card\,red}C'_0\leq\pi_1^{\pi_0}$.
% element: margin note | left lower middle, slanted | Cette majoration
% slanted
\margin{Cette majoration n'est pas prouvée, et \unsure{pro.b.blement|prob.blement|probablement} \unsure{fausse} \unsure{si on ne fait pas} des hypothèses un peu plus \unsure{strictes}}

% element: text | lower third, corollary heading and statement | Corollaire 4.6.1.2
\deleted{\unsure{3°)|3)}} \uline{Corollaire} 4.6.1.2. $M$ et $C_0$ étant comme dans 2°) ci-dessus, soit $\unsure{C_0|C}$ \unsure{le|la} plus petite sous-catégorie str. pleine de $M$ \deleted{telle que} contenant $C_0$, et stable par $\unsure{\underline{\lim}|\varinjlim}$ de cardinal $\leq\pi_0$ (celles qui sont repr. dans $M$). Si $C_0$ est équivalente à une petite catégorie, il en est de $\unsure{=|\bar m}$ de $C$. Plus précisément, si $\pi_1$ est un cardinal infini, \inserted{et $\operatorname{card\,red}C_0\leq\unsure{\pi_1|\pi}$,} alors

\scanpage{227}
\mspage{143}
% archive sheet number: 226
% element: text | top, corollary conclusion | alors on a
alors on a $\operatorname{card\,réd}C\leq\pi_1^{\pi_0}$.
% element: margin note | upper left, slanted | majoration non prouvee
% slanted
\margin{majoration \unsure{non prouvée}, \unsure{prob.blemt|prob.blement|probablement} fausse, cf \unsure{annotation} ci-dessus.}

% element: text | upper third, proof | Le 3 se prouve
\unsure{Le|La} 3°) se prouve par la démonstration \deleted{\unsure{préc.|puis}} transfinie précédente, \deleted{\unsure{on|les}} on prouve par récurrence transfinie que
% element: display | upper third, inductive bound | card red C alpha
\[
\operatorname{card\,réd}C_\alpha\leq\pi_1^{\pi_0}
\]
% element: text | upper middle, initial and successor cases | pour tt alpha
pour \unsure{tt|$\forall$} $\alpha$. C'est \unsure{o.K|o.k} pour $\alpha=0$ par hyp (puisque $\pi_1\leq\pi_1^{\pi_0}$), \deleted{c'est vrai} si c'est vrai pour $C_\alpha$, alors 2°) assure que
% element: display | middle, successor estimate | card red C alpha power estimate
\[
\operatorname{card\,red}C_\alpha\leq(\pi_1^{\pi_0})^{\pi_0}=\pi_1^{\pi_0\cdot\pi_0}=\pi_1^{\pi_0},
\]
% element: text | middle, limit case | enfin si alpha est un ordinal limite
enfin, si $\alpha$ est un ordinal limite $<\alpha_0$, on a $C_\alpha=\bigcup_{\alpha'<\alpha}C_{\alpha'}$, donc
% element: display | middle lower, limit estimate | card red C alpha product estimate
\[
\operatorname{card\,réd}C_\alpha\leq\pi_0\cdot\pi_1^{\pi_0}=\pi_1^{\pi_0}
\]
% element: text | lower middle, justification and successor cardinal | vu que
(\unsure{vu|car} que $\pi_1^{\pi_0}\geq\pi_0$), \deleted{Donc \unsure{on a|il est}} \inserted{en vertu de} \deleted{par} 1° ci-dessus, Et c'est dans ce cas que $\operatorname{card}I_\alpha\leq\pi_0$. Soit $\pi'_0$ le cardinal \inserted{successeur} de $\pi_0$, alors \unsure{la|le} $\unsure{=|\bar m}$ majoration nous donne
% element: display | lower third, bound at alpha zero | card red C alpha zero
\[
\operatorname{card\,red}C_{\alpha_0}\leq\pi'_0\pi_1^{\pi_0}.
\]
% element: text | lower third, cardinal comparison | Mais
Mais \unsure{on a|...} $\pi_1^{\pi_0}\geq 2^{\pi_0}>\pi_0$, donc $\pi_1^{\pi_0}\geq\pi'_0$, donc $\pi'_0\pi_1^{\pi_0}=\pi_1^{\pi_0}$, d'où
% element: display | foot, final estimate | card red C alpha zero final bound
\[
\operatorname{card\,red}C_{\alpha_0}\leq\pi_1^{\pi_0},
\]
% element: text | foot, conclusion | ce qui est
ce qui est la majoration annoncée.

\scanpage{228}
\mspage{144}
% archive sheet number: 227
% element: text | top, second conclusion of proposition | Prouvons 4.6.1 2
Prouvons 4.6.1. 2°, i.e $C=M_\pi$. On a $C\subset M_{\pi_0}$, puisque par hypothèse $S\subset M_{\pi_0}$, et on sait que $M_{\pi_0}$ est stable par $\unsure{\underline{\lim}|\varinjlim}$ de cardinal $\leq\pi_0$. D'autre part: $M_{\pi_0}\subset C$, en vertu de 4.5.14 $a\Rightarrow c$ appliqué à $C$, qui est strictement génératrice dans $M$, $S$ l'étant et $C\supset S$, et stable par $\underline{\lim}$ de cardinal $\leq\pi_0$ par construction.

% element: text | upper middle, transition to third conclusion | Prouvons enfin
\uline{Prouvons enfin} 3°. Cela \deleted{\unsure{va résulter du|va résulter des}} \inserted{\unsure{équivaut} à ceci:}

% element: text | middle, lemma and revised assumptions | Lemme 4.6.1.3
\uline{Lemme} 4.6.1.3. Soit $M$ une $\unsure{U|\mathcal U|\mathscr U}$-catégorie, \inserted{stable par \uline{petites} $\underline{\lim}$,} $\pi_0$ un cardinal, $M_0$ une sous-catégorie de $M$ \deleted{str. pleine, \unsure{str. pleine,} \unsure{équiv. à une} petite catégorie,} \deleted{telle que $M_0$ soit} \inserted{\uline{\unsure{stricte|strictement} génératrice},} \deleted{\unsure{et}} stable par $\varinjlim$ de cardinal $\leq\pi_0$, \deleted{Alors} \unsure{\inserted{et formée d'objets $\pi_0$-accessibles.}|\inserted{et} formée d'objets $\pi_0$-accessibles.} \deleted{\unsure{Partons du cas où|Prouvons dans ce cas}} (En vertu de 4.6.1.\unsure{2|4}, 2° déjà prouvé, appliqué à $\deleted{\unsure{C}}S=M_0$, \inserted{si $M_0$ est str. pleine,} \unsure{on a} $M_0=M_{\pi_0}$) Alors le foncteur canonique
% element: diagram | lower middle, canonical functor | Ind prime pi zero M zero line M
\diagram{\begin{tikzcd}[ampersand replacement=\&]
\mathrm{Ind}'_{\pi_0}(M_0)\arrow[r,dash]\&M
\end{tikzcd}}
% element: text | lower third, conclusion and full faithfulness | est une equivalence
est une équivalence de catégories.

Le fait que ce foncteur est pl. fidèle résulte aussitôt de l'hypothèse que les objets de $M_0$ sont $\pi_0$-accessibles:
% element: display | foot, Hom comparison and vertical membership annotations | Hom M colimits
% The tiny signs below x_i and y_j are membership signs rotated clockwise, with M_0 below them.
\[
\begin{aligned}
\mathrm{Hom}_M\biggl(\varinjlim_{\unsure{j|i}\in I}\underset{\substack{\tikz[baseline=-0.4ex]{\node[rotate=-90,inner sep=0pt] {$\scriptstyle\in$};}\\M_0}}{x_i},\varinjlim_{j\in J}\underset{\substack{\tikz[baseline=-0.4ex]{\node[rotate=-90,inner sep=0pt] {$\scriptstyle\in$};}\\M_0}}{y_j}\biggr)
&\simeq\varprojlim_{i\in I}\underbrace{\mathrm{Hom}(x_i;\varinjlim_{j\in J}y_j)}_{\displaystyle\varinjlim_{j\in J}\mathrm{Hom}(x_i;y_j)}\\
&\simeq\mathrm{Hom}_{\mathrm{Ind}_{\pi}(M)}((x_i),(y_j))\unsure{,|.}
\end{aligned}
\]
